% ============================================================================
\section{Paradigm 3: Logic Gate Networks}
\label{sec:lgn}
% ============================================================================

% --- Intuition ---
\subsection{The Idea in Plain English}

Binary Neural Networks hardcode the operation: it is always XNOR.
Tsetlin Machines hardcode the operation: it is always AND (within
clauses). But what if the network could \emph{choose} which logic
gate to use at each connection?

A Logic Gate Network~\cite{petersen2022difflogic} is a layered network where each ``neuron'' is
a 2-input logic gate, and both the gate's \emph{inputs} (which two
signals to connect) and its \emph{type} (AND, OR, XOR, NAND, etc.)
are learned during training.

\begin{analogy}
Imagine building a circuit on a breadboard. You have a box of
different logic gate chips (AND, OR, XOR, NAND, NOR, XNOR). A
Logic Gate Network is like having a robot that tries different chips
in each socket and different wiring configurations, testing each
combination, and gradually settling on the circuit that best solves
the classification problem.
\end{analogy}

% --- Mechanism ---
\subsection{How It Works}

\subsubsection{The 16 Possible 2-Input Gates}

Two binary inputs can produce $2^{2^2} = 16$ possible logic functions.
These include the familiar gates plus some degenerate cases:

\begin{figure}[htbp]
\centering
\begin{tikzpicture}[
  cell/.style={minimum width=1.0cm, minimum height=0.5cm,
    font=\scriptsize\ttfamily, inner sep=0pt},
  header/.style={cell, font=\scriptsize\sffamily\bfseries},
  >=Stealth
]
  % Table header
  \node[header] at (0, 0.8) {$a$};
  \node[header] at (1, 0.8) {$b$};
  \draw[thick] (-0.5, 0.5) -- (10.5, 0.5);

  % Input rows
  \foreach \a/\b/\row in {0/0/0, 0/1/-0.5, 1/0/-1, 1/1/-1.5} {
    \node[cell] at (0, \row) {\a};
    \node[cell] at (1, \row) {\b};
  }

  % Selected gates - explicit nodes to avoid nested foreach issues
  % AND
  \node[header] at (2.5, 0.8) {AND};
  \node[cell] at (2.5, 0) {0}; \node[cell] at (2.5, -0.5) {0};
  \node[cell] at (2.5, -1) {0}; \node[cell, fill=bitblue!15] at (2.5, -1.5) {1};
  % OR
  \node[header] at (3.7, 0.8) {OR};
  \node[cell] at (3.7, 0) {0}; \node[cell, fill=bitblue!15] at (3.7, -0.5) {1};
  \node[cell, fill=bitblue!15] at (3.7, -1) {1}; \node[cell, fill=bitblue!15] at (3.7, -1.5) {1};
  % XOR
  \node[header] at (4.9, 0.8) {XOR};
  \node[cell] at (4.9, 0) {0}; \node[cell, fill=bitblue!15] at (4.9, -0.5) {1};
  \node[cell, fill=bitblue!15] at (4.9, -1) {1}; \node[cell] at (4.9, -1.5) {0};
  % NAND
  \node[header] at (6.1, 0.8) {NAND};
  \node[cell, fill=bitblue!15] at (6.1, 0) {1}; \node[cell, fill=bitblue!15] at (6.1, -0.5) {1};
  \node[cell, fill=bitblue!15] at (6.1, -1) {1}; \node[cell] at (6.1, -1.5) {0};
  % NOR
  \node[header] at (7.3, 0.8) {NOR};
  \node[cell, fill=bitblue!15] at (7.3, 0) {1}; \node[cell] at (7.3, -0.5) {0};
  \node[cell] at (7.3, -1) {0}; \node[cell] at (7.3, -1.5) {0};
  % XNOR
  \node[header] at (8.5, 0.8) {XNOR};
  \node[cell, fill=bitblue!15] at (8.5, 0) {1}; \node[cell] at (8.5, -0.5) {0};
  \node[cell] at (8.5, -1) {0}; \node[cell, fill=bitblue!15] at (8.5, -1.5) {1};

  % Separator
  \draw[bitgray, thick] (1.5, 1.1) -- (1.5, -1.8);
\end{tikzpicture}
\caption{Truth tables for six commonly used 2-input gates. There are 16
  possible functions in total. An LGN learns which function each node
  should implement.}
\label{fig:truth-tables}
\end{figure}

\subsubsection{Differentiable Gate Selection}

The core trick of LGNs is making the discrete choice of gate type
\emph{differentiable} so that gradient descent can be used for
training.

Each node maintains a probability distribution over the 16 possible
gate functions. During training, the output is a weighted average:

\begin{equation}
  \text{output} = \sum_{g=0}^{15} p_g \cdot f_g(a, b)
  \label{eq:lgn-soft}
\end{equation}

where $p_g$ is the (learnable) probability of gate $g$ and $f_g(a,b)$
is gate $g$'s output for inputs $a, b$.

\begin{figure}[htbp]
\centering
\begin{tikzpicture}[
  node/.style={draw, circle, minimum size=0.8cm, font=\tiny\sffamily,
    inner sep=0pt, line width=0.8pt},
  >=Stealth
]
  % Input layer
  \foreach \i in {1,...,6} {
    \node[node, fill=lightbg] (in\i) at (0, 5.5-\i) {$x_\i$};
  }

  % Hidden layer 1
  \foreach \i/\gate/\col in {1/AND/bitblue, 2/XOR/bitorange,
    3/OR/bitgreen, 4/NAND/bitpurple} {
    \node[node, fill=\col!15] (h1\i) at (3, 5-\i) {\gate};
  }

  % Hidden layer 2
  \foreach \i/\gate/\col in {1/XOR/bitorange, 2/AND/bitblue,
    3/OR/bitgreen} {
    \node[node, fill=\col!15] (h2\i) at (6, 4.5-\i) {\gate};
  }

  % Output
  \node[node, fill=bitblue!15] (out) at (9, 3) {AND};

  % Connections (selected, not all-to-all)
  \draw[->, thick] (in1) -- (h11);
  \draw[->, thick] (in2) -- (h11);
  \draw[->, thick] (in1) -- (h12);
  \draw[->, thick] (in3) -- (h12);
  \draw[->, thick] (in4) -- (h13);
  \draw[->, thick] (in5) -- (h13);
  \draw[->, thick] (in3) -- (h14);
  \draw[->, thick] (in6) -- (h14);

  \draw[->, thick] (h11) -- (h21);
  \draw[->, thick] (h12) -- (h21);
  \draw[->, thick] (h12) -- (h22);
  \draw[->, thick] (h13) -- (h22);
  \draw[->, thick] (h13) -- (h23);
  \draw[->, thick] (h14) -- (h23);

  \draw[->, thick] (h21) -- (out);
  \draw[->, thick] (h23) -- (out);

  % Labels
  \node[font=\small\sffamily, above] at (0, 5) {Inputs};
  \node[font=\small\sffamily, above] at (3, 5) {Layer 1};
  \node[font=\small\sffamily, above] at (6, 5) {Layer 2};
  \node[font=\small\sffamily, above] at (9, 4) {Output};

  % Annotation
  \node[font=\scriptsize\sffamily, bitgray, align=center] at (4.5, -0.5)
    {Gate types are learned; connectivity is random\\but \emph{effectively}
    learned (see text)};
\end{tikzpicture}
\caption{A Logic Gate Network. Each node is a 2-input logic gate whose
  type is learned through differentiable relaxation. The connectivity
  is initially random but effectively learned through gate-type
  selection (see below). At inference time, this is a pure logic
  circuit with zero arithmetic.}
\label{fig:lgn-architecture}
\end{figure}

\subsubsection{How Connectivity Is ``Learned''}
\label{sec:lgn-connectivity}

The equations in the next subsection show how gate \emph{types} are
learned, but how does the network learn its \emph{wiring}---which two
signals feed into each gate?

In the original DiffLogic architecture~\cite{petersen2022difflogic},
the answer is surprisingly indirect: \textbf{connections are assigned
randomly and remain fixed.} Each gate is randomly wired to two outputs
from the previous layer at initialisation and never rewired.

This sounds limiting, but it works because of a key observation:
among the 16 possible 2-input Boolean functions, six of them
\emph{effectively ignore one or both inputs}:

\begin{table}[htbp]
\centering
\caption{Of the 16 possible 2-input Boolean functions, six ignore
  at least one input. By learning to become one of these functions,
  a gate can effectively ``disconnect'' an unwanted input.}
\label{tab:gate-disconnect}
\begin{tabular}{@{}llll@{}}
  \toprule
  \textbf{Gate function} & \textbf{Output} & \textbf{Ignores}
    & \textbf{Effect} \\
  \midrule
  Constant 0  & always 0             & both $a$ and $b$
    & dead gate (pruned) \\
  Constant 1  & always 1             & both $a$ and $b$
    & dead gate (pruned) \\
  Pass-through $a$  & $a$            & input $b$
    & wire from $a$ only \\
  Pass-through $b$  & $b$            & input $a$
    & wire from $b$ only \\
  NOT $a$     & $\overline{a}$       & input $b$
    & inverter on $a$ \\
  NOT $b$     & $\overline{b}$       & input $a$
    & inverter on $b$ \\
  \bottomrule
\end{tabular}
\end{table}

\begin{keyinsight}
Connectivity learning is \emph{subsumed by gate-type learning.} When
the training process selects a gate type, it simultaneously decides
the effective wiring:
\begin{itemize}[leftmargin=1em]
  \item A gate that learns ``AND'' uses both inputs.
  \item A gate that learns ``pass-through $a$'' effectively
    disconnects input $b$, acting as a simple wire.
  \item A gate that learns ``constant 0'' effectively disconnects
    from the network entirely.
\end{itemize}
The network is deliberately \emph{over-provisioned}---it starts with
many more gates than needed, and training prunes unnecessary
connections by turning gates into pass-throughs or constants.
This is why the single softmax-over-16-gates mechanism
(\cref{eq:lgn-soft}) is sufficient to learn both gate types
\emph{and} effective connectivity.
\end{keyinsight}

\subsubsection{Training: From Soft to Hard}

During training, the network operates in ``soft'' mode: gate outputs
are real-valued probabilities, and gradients flow through the weighted
sum. As training progresses, the distributions sharpen until each node
has a clear winner. Two common techniques drive this sharpening:
\emph{temperature annealing} gradually lowers a temperature parameter
that controls how peaked the softmax distribution is---high temperature
gives near-uniform probabilities, low temperature approaches a hard
one-hot selection. The \emph{Gumbel-softmax} trick adds carefully
calibrated random noise before the softmax, which lets the network
explore different gate choices early in training while still permitting
gradient flow (unlike a hard \texttt{argmax}, which is not
differentiable).

At inference time, each node commits to its most probable gate type
and operates in pure binary logic---no floating-point, no
probabilities.

\begin{figure}[htbp]
\centering
\begin{tikzpicture}[>=Stealth]
  % Training phase
  \begin{scope}
    \node[font=\small\sffamily\bfseries] at (2, 3) {During Training};
    \draw[fill=bitblue!20, draw=bitblue] (0, 0) rectangle (0.6, 1.5);
    \draw[fill=bitorange!40, draw=bitorange] (0.8, 0) rectangle (1.4, 2.2);
    \draw[fill=bitgreen!20, draw=bitgreen] (1.6, 0) rectangle (2.2, 0.8);
    \draw[fill=bitpurple!15, draw=bitpurple] (2.4, 0) rectangle (3.0, 0.3);
    \draw[fill=bitgray!10, draw=bitgray] (3.2, 0) rectangle (3.8, 0.2);
    \node[font=\tiny\sffamily, below, rotate=45] at (0.3, -0.1) {AND};
    \node[font=\tiny\sffamily, below, rotate=45] at (1.1, -0.1) {XOR};
    \node[font=\tiny\sffamily, below, rotate=45] at (1.9, -0.1) {OR};
    \node[font=\tiny\sffamily, below, rotate=45] at (2.7, -0.1) {NAND};
    \node[font=\tiny\sffamily, below, rotate=45] at (3.5, -0.1) {\ldots};
    \node[font=\scriptsize\sffamily, bitgray] at (2, -1.2)
      {Soft probabilities};
  \end{scope}

  % Arrow
  \draw[->, very thick, bitgray] (5, 1) -- (6.5, 1)
    node[midway, above, font=\scriptsize\sffamily] {anneal};

  % Inference phase
  \begin{scope}[xshift=7.5cm]
    \node[font=\small\sffamily\bfseries] at (2, 3) {At Inference};
    \draw[fill=bitblue!5, draw=bitgray!30] (0, 0) rectangle (0.6, 0.05);
    \draw[fill=bitorange!60, draw=bitorange, line width=1.5pt]
      (0.8, 0) rectangle (1.4, 2.5);
    \draw[fill=bitgreen!5, draw=bitgray!30] (1.6, 0) rectangle (2.2, 0.05);
    \draw[fill=bitpurple!5, draw=bitgray!30] (2.4, 0) rectangle (3.0, 0.05);
    \draw[fill=bitgray!5, draw=bitgray!30] (3.2, 0) rectangle (3.8, 0.05);
    \node[font=\tiny\sffamily, below, rotate=45] at (0.3, -0.1) {AND};
    \node[font=\tiny\sffamily\bfseries, below, rotate=45] at (1.1, -0.1) {XOR};
    \node[font=\tiny\sffamily, below, rotate=45] at (1.9, -0.1) {OR};
    \node[font=\tiny\sffamily, below, rotate=45] at (2.7, -0.1) {NAND};
    \node[font=\tiny\sffamily, below, rotate=45] at (3.5, -0.1) {\ldots};
    \node[font=\scriptsize\sffamily, bitgray] at (2, -1.2)
      {Hard selection: XOR};
  \end{scope}
\end{tikzpicture}
\caption{Gate selection during training (soft, differentiable
  probabilities) versus inference (hard, discrete gate commitment).
  Temperature annealing gradually sharpens the distribution.}
\label{fig:lgn-annealing}
\end{figure}

% --- Formalisation ---
\subsection{Mathematical Summary}
\label{sec:lgn-math}

Let each node $i$ in layer $\ell$ receive two binary inputs
$a_i^{(\ell)}, b_i^{(\ell)} \in \{0, 1\}$ from the previous layer.

The node maintains a logit vector $\bm{\alpha}_i \in \mathbb{R}^{16}$,
converted to probabilities via softmax with temperature $\tau$:
\begin{equation}
  p_g^{(i)} = \frac{\exp(\alpha_g^{(i)} / \tau)}
    {\sum_{g'=0}^{15} \exp(\alpha_{g'}^{(i)} / \tau)}
  \label{eq:lgn-softmax}
\end{equation}

The soft output during training is:
\begin{equation}
  o_i^{(\ell)} = \sum_{g=0}^{15} p_g^{(i)} \cdot
    \operatorname{TruthTable}(g, a_i^{(\ell)}, b_i^{(\ell)})
  \label{eq:lgn-soft-output}
\end{equation}

where $\operatorname{TruthTable}(g, a, b) \in \{0, 1\}$ is the output
of the $g$-th Boolean function for inputs $(a, b)$.

At inference ($\tau \to 0$):
\begin{equation}
  o_i^{(\ell)} = \operatorname{TruthTable}\!\left(
    \arg\max_g \alpha_g^{(i)},\; a_i^{(\ell)},\; b_i^{(\ell)}\right)
  \label{eq:lgn-hard-output}
\end{equation}

\subsection{Performance and Scale}

LGNs are relatively new. Published results include:

\begin{table}[htbp]
\centering
\caption{Logic Gate Network results on MNIST. Connection optimisation
  substantially changes the gate-count picture, but the newest result is a
  preprint and needs independent replication~\cite{petersen2022difflogic,mommen2026fullytrainable}.}
\label{tab:lgn-perf}
\begin{tabular}{@{}lrrl@{}}
  \toprule
  \textbf{Network} & \textbf{MNIST} & \textbf{Gates} & \textbf{Note} \\
  \midrule
  DiffLogic (small)~\cite{petersen2022difflogic} & 97.7\% & 48k  & 6 layers, 8k gates/layer \\
  DiffLogic~\cite{petersen2022difflogic}          & 98.5\% & 384k & 6 layers, 64k gates/layer \\
  Connection-optimised LGN~\cite{mommen2026fullytrainable} & 98.92\% & 16k & 2 layers, 8k gates/layer \\
  \bottomrule
\end{tabular}
\end{table}

Recent work changes two important limitations of the original LGN
formulation. A reparameterisation reduces binary-input model size by
$4\times$, accelerates the backward pass, and reports stable CIFAR-100
accuracy~\cite{ruttgers2025light}. Separately, learning the connections as
well as the gate types reports 98.92\% MNIST accuracy with two layers of
8{,}000 gates---almost $50\times$ fewer gates than a comparable
fixed-connection construction~\cite{mommen2026fullytrainable}. These
results make sparse, learned Boolean composition a stronger CPU research
lead, but they do not yet establish transformer-scale circuits.

\subsection{Strengths and Limitations}

\begin{description}[leftmargin=1em]
  \item[Strengths:]
    Pure logic at inference, learned gate type and connectivity, and a
    plausible path to sparse, branch-light CPU evaluation when gate outputs
    are stored and traversed in contiguous packed arrays.

  \item[Limitations:]
    2-input restriction can require deep networks; training is slower than
    standard backpropagation (16-way softmax per node); and a gate graph
    with irregular fan-out can lose to a small SIMD integer MLP on CPU
    despite performing no multiplications.
\end{description}
